% =============================================================================
% chapter4_cuts.tex
%
% Five slimmed replacement paragraphs for chapter4.tex. Each block below
% replaces a specific paragraph (or paragraph group) in chapter4.tex and
% trades framework/implementation minutiae for tighter exposition while
% preserving every architectural idea.
%
% This file is NOT \input by main.tex. It is a staging file: copy each block
% over the matching range in chapter4.tex when ready.
%
% Cut 1 — replaces line 246          (§4.3.2 JWT request handling)
% Cut 2 — replaces lines 442--450    (§4.4.5 Frontend Authentication)
% Cut 3 — replaces line 464          (§4.4.6 Login/Register page)
% Cut 4 — replaces lines 490--498    (§4.4.6 Results page)
% Cut 5 — replaces lines 516--530    (§4.4.7 Reusable Components)
% =============================================================================


% -----------------------------------------------------------------------------
% Cut 1 — replaces chapter4.tex line 246
% Section: §4.3.2 Authentication and Security
% Before:  one long paragraph describing OncePerRequestFilter, CustomUserDetailsService,
%          and UserPrincipal field contents.
% After:   same paragraph minus filter-class name and UserPrincipal field list.
% -----------------------------------------------------------------------------
On every incoming request, the \texttt{JwtAuthenticationFilter} extracts the token from the \texttt{Authorization: Bearer} header, validates it, loads the corresponding user and populates the Spring \texttt{SecurityContext}. If no valid token is present, the \texttt{JwtAuthenticationEntryPoint} returns a \texttt{401 Unauthorized} response. Passwords are hashed using BCrypt \cite{Provos1999}, and the authenticated user's identity is injected into controller methods via the \texttt{@AuthenticationPrincipal} annotation.


% -----------------------------------------------------------------------------
% Cut 2 — replaces chapter4.tex lines 442--450
% Section: §4.4.5 Authentication Integration
% Before:  five paragraphs detailing AuthService (BehaviorSubject, localStorage keys,
%          login/register/logout/isAuthenticated methods), authInterceptor
%          (HttpInterceptorFn, inject(), bypass for /auth/ endpoints), errorInterceptor
%          (RxJS catchError, 401 handling), and authGuard (CanActivateFn, route list).
% After:   one paragraph naming the four pieces and what each does.
% -----------------------------------------------------------------------------
\par The frontend authentication layer consists of four cooperating pieces. The \texttt{AuthService} singleton manages JWT storage in the browser's \texttt{localStorage} and exposes the current user as a reactive observable consumed by the navigation bar. An \texttt{authInterceptor} automatically attaches the \texttt{Authorization: Bearer} header to every outgoing API request (except public authentication endpoints), while an \texttt{errorInterceptor} catches \texttt{401} responses and triggers a logout with redirect to the login page. Finally, a functional \texttt{authGuard} protects the \texttt{/upload}, \texttt{/results/:jobId}, and \texttt{/history} routes by checking the token's presence before allowing navigation.


% -----------------------------------------------------------------------------
% Cut 3 — replaces chapter4.tex line 464
% Section: §4.4.6 Pages (Login/Register)
% Before:  detailed walkthrough of FormsModule, ngModel, canSubmit getter,
%          AuthService.login() flow, 6-char min password, navigation links.
% After:   one sentence covering credential forms with inline error feedback.
% -----------------------------------------------------------------------------
The \texttt{LoginPageComponent} and \texttt{RegisterPageComponent} provide standard credential forms with inline error feedback; both navigate to the upload page on success.


% -----------------------------------------------------------------------------
% Cut 4 — replaces chapter4.tex lines 490--498
% Section: §4.4.6 Pages (Results)
% Before:  five paragraphs covering route param read, getJobResults call chain,
%          header layout with Re-Analyze dropdown + Export JSON, modal dialog
%          details, diff banner conditional, summary card grid, breakdown,
%          two-column anti-pattern list + dependency graph layout.
% After:   one paragraph naming every component and feature, with Figure ref.
% -----------------------------------------------------------------------------
The \texttt{ResultsPageComponent} (Figure~\ref{fig:results-page}) is the primary output view. It displays a health score gauge, summary metric cards, a per-category score breakdown, an expandable list of detected anti-patterns with code evidence, and an interactive dependency graph rendered with Cytoscape.js. A re-analysis dropdown allows uploading a new ZIP or cloning a different URL, and an export button downloads the full result as JSON. When the analysis is a re-analysis, a diff banner compares the current and previous results.


% -----------------------------------------------------------------------------
% Cut 5 — replaces chapter4.tex lines 516--530
% Section: §4.4.7 Reusable Components
% Before:  eight paragraphs (one intro + one per component) with deep
%          implementation detail: SVG circumference math, three-tier color
%          coding, eleven-card diff banner grid, CoSE layout parameters,
%          drag DOM events, file size formatting, language detection.
% After:   one paragraph listing all seven components with their roles.
%
% Optional restore: if you want to keep the DiffBanner's "inverse" mode
% (a non-obvious UX choice for metrics where decrease = improvement),
% append the following sentence to the paragraph below:
%   The \texttt{DiffBannerComponent} additionally supports an \texttt{inverse}
%   colour-coding mode for metrics where a decrease is an improvement
%   (e.g.\ fewer anti-patterns).
% -----------------------------------------------------------------------------
\par The frontend defines seven standalone reusable components composed into pages via Angular's input and output bindings: \texttt{HealthScoreComponent} (circular SVG gauge with letter grade), \texttt{Health\-Score\-Break\-down\-Com\-po\-nent} (per-category collapsible score cards), \texttt{Diff\-Ban\-ner\-Com\-po\-nent} (metric deltas between consecutive analyses), \texttt{De\-pen\-den\-cy\-Graph\-Com\-po\-nent} (interactive Cytoscape.js graph with CoSE layout), \texttt{File\-Up\-load\-Com\-po\-nent} (drag-and-drop ZIP selection), \texttt{Code\-Snip\-pet\-View\-er\-Com\-po\-nent} (inline source evidence with line highlighting), and \texttt{Pro\-gress\-Track\-er\-Com\-po\-nent} (multi-step pipeline progress indicator). Each component is purely presentational and receives its data via Angular's \texttt{@Input} bindings from the parent page.
